\section*{Chapter \ref{ch:g6:inside-the-earth} --- Inside the Earth}

\begin{solution}{exo:g6:inside-the-earth:1}
The \omterm{def:g6:inside-the-earth:layers}{crust}, the \omterm{def:g6:inside-the-earth:layers}{mantle}, the \omterm{def:g6:inside-the-earth:layers}{outer core}, the \omterm{def:g6:inside-the-earth:layers}{inner core}. The \omterm{def:g6:inside-the-earth:layers}{outer core} is
liquid. The outer and \omterm{def:g6:inside-the-earth:layers}{inner core} are made of metal (mostly iron, with some
nickel); the \omterm{def:g6:inside-the-earth:layers}{crust} and the \omterm{def:g6:inside-the-earth:layers}{mantle} are made of \omterm{def:g3:stones-pebbles-sand:rock}{rock}.
\end{solution}

\begin{solution}{exo:g6:inside-the-earth:2}
(a) False: the \omterm{def:g6:inside-the-earth:layers}{crust} is thicker under the \omterm{def:g3:continents-and-oceans:continent}{continents} (about 30~km) than
under the \omterm{def:g2:water-on-earth:ocean}{oceans} (about 5~km). (b) False: the \omterm{def:g6:inside-the-earth:layers}{core} is made of metal, an
alloy of iron and nickel. (c) True. (d) True: 12\,262~m, about 12.3~km.
\end{solution}

\begin{solution}{exo:g6:inside-the-earth:3}
(a) At 2891~km, about 2900~km. (b) At 5150~km. (c) $5150 - 2891 = 2259$:
about 2260~km.
\end{solution}

\begin{solution}{exo:g6:inside-the-earth:4}
``The shell is like the \underline{\omterm{def:g6:inside-the-earth:layers}{crust}}, the white is like the
\underline{\omterm{def:g6:inside-the-earth:layers}{mantle}} and the yolk is like the \underline{\omterm{def:g6:inside-the-earth:layers}{core}}. But unlike the
yolk, the \underline{\omterm{def:g6:inside-the-earth:layers}{core}} has two parts: a liquid \underline{\omterm{def:g6:inside-the-earth:layers}{outer core}}
and a solid \underline{\omterm{def:g6:inside-the-earth:layers}{inner core}}.''
\end{solution}

\begin{solution}{exo:g6:inside-the-earth:5}
$8.849 \div 6371 \approx 0.00139$, and $0.00139 \times 100 \approx 0.14$:
Everest is only about \textbf{0.14~\%} of the radius. Seen from far away,
the Earth is a very smooth ball: its highest \omterm{def:g1:landscapes-around-us:landform}{mountain} is a tiny bump.
\end{solution}

\begin{solution}{exo:g6:inside-the-earth:6}
(a) $30 \div 5 = 6$: six times thicker. (b) $100 \div 5 = 20$: twenty times
thicker.
\end{solution}

\begin{solution}{exo:g6:inside-the-earth:7}
(a) $12 + 1 \times 25 = 37~^\circ$C. (b) $12 + 1.5 \times 25 = 12 + 37.5 =
49.5~^\circ$C. (c) $12 + 2 \times 25 = 62~^\circ$C.
\end{solution}

\begin{solution}{exo:g6:inside-the-earth:8}
The rise is $74 - 14 = 60$ degrees over 2~km, so $60 \div 2 =
\textbf{30}$~degrees per kilometre. Yes: 30 is between the eastern average
(25) and the western average (34).
\end{solution}

\begin{solution}{exo:g6:inside-the-earth:9}
(a) Just above the \omterm{def:g6:inside-the-earth:layers}{core}, about 13.7~km per second; just below its top,
about 8.1~km per second: the P waves suddenly slow down. (b) Between
2891~km and 5150~km, in the \omterm{def:g6:inside-the-earth:layers}{outer core}. (c) S waves cannot travel through a
liquid: the \omterm{def:g6:inside-the-earth:layers}{outer core} is liquid.
\end{solution}

\begin{solution}{exo:g6:inside-the-earth:10}
(a) $12.3 \div 30 = 0.41$: about \textbf{41~\%} of an average continental
\omterm{def:g6:inside-the-earth:layers}{crust}. (b) No: under a \omterm{def:g3:continents-and-oceans:continent}{continent} the \omterm{def:g6:inside-the-earth:layers}{mantle} begins below the \omterm{def:g6:inside-the-earth:layers}{crust}, about
30~km down on average; the borehole stopped far above it, in the \omterm{def:g6:inside-the-earth:layers}{crust}.
\end{solution}

\begin{solution}{exo:g6:inside-the-earth:11}
(a) $29 \times 6371 = 184\,759$, and $15 + 184\,759 = 184\,774$: about
185\,000~degrees. (b) The laboratory gives about 6000~degrees at the edge of
the \omterm{def:g6:inside-the-earth:layers}{inner core}: Leo's number is about 30~times too large
($184\,774 \div 6000 \approx 31$). (c) The rate of 29~degrees per kilometre
was measured only in the first 2~km or so. Leo uses it thousands of
kilometres beyond the measurements; deep down the temperature rises much
more slowly. A rate may only be used where it was measured.
\end{solution}

\begin{solution}{exo:g6:inside-the-earth:12}
(a) If the \omterm{def:g6:inside-the-earth:layers}{crust} is lighter than the average, the material deeper down must
be denser than 5.5, to bring the average up. (b) In the \omterm{def:g6:inside-the-earth:layers}{outer core} and the
\omterm{def:g6:inside-the-earth:layers}{inner core} (from about 9.9 to about 13~grams per cubic centimetre); the
\omterm{def:g6:inside-the-earth:layers}{mantle} only just reaches 5.6 at its very bottom. (c) Denser materials sink
below lighter ones, as a stone sinks in water: the heavy metal ended up in
the middle and the lighter \omterm{def:g3:stones-pebbles-sand:rock}{rocks} on top.
\end{solution}

\begin{solution}{pb:g6:inside-the-earth:1}
\textbf{1.} $2 \times 6371 = 12\,742$~km. The model is 2~m across, that is
2000~mm.

\textbf{2.} $12\,742 \div 2000 \approx 6.37$: 1~mm of the model stands for
about 6.37~km of the real Earth.

\textbf{3.} $8.849 \div 6.371 \approx 1.39$: Everest is about
\textbf{1.4~mm} high on the model, a grain of \omterm{def:g3:stones-pebbles-sand:sand}{sand}.

\textbf{4.} $12.3 \div 6.371 \approx 1.9$~mm. The centre of the model is
1000~mm (1~m) below the surface: the borehole is about 520~times shorter
($1000 \div 1.9 \approx 520$).

\textbf{5.} $2891 \div 6.371 \approx 453.8$: the top of the \omterm{def:g6:inside-the-earth:layers}{outer core} is
painted about \textbf{454~mm} below the surface.

\textbf{6.} $5150 \div 6.371 \approx 808.3$: about 808~mm below the
surface. The \omterm{def:g6:inside-the-earth:layers}{inner core}'s radius on the model is $1000 - 808 = 192$~mm.

\textbf{7.} The \omterm{def:g6:inside-the-earth:layers}{mantle}: from about 5~mm to 454~mm, about 449~mm. The \omterm{def:g6:inside-the-earth:layers}{outer
core} takes $808 - 454 = 354$~mm and the \omterm{def:g6:inside-the-earth:layers}{inner core} 192~mm.

\textbf{8.} $30 \div 6.371 \approx 4.7$~mm: a thin band, which can be
painted with a fine brush.

\textbf{9.} $15 + 2 \times 29 = 73~^\circ$C. On the model, 2~km is
$2 \div 6.371 \approx 0.31$~mm below the surface: hardly the thickness of
three sheets of paper.

\textbf{10.} $4 \div 6.371 \approx 0.63$~mm, against about 1.9~mm for the
borehole: the mine is shallower, by about $1.9 - 0.6 = 1.3$~mm.

\textbf{11.} Whether a substance melts depends on the temperature and on
the pressure. The \omterm{def:g6:inside-the-earth:layers}{inner core} is squeezed hardest of all, by about
3.3~million times the pressure of the air: squeezed iron melts only at a
higher temperature, so the iron of the \omterm{def:g6:inside-the-earth:layers}{inner core} stays solid although it
is hotter than the \omterm{def:g6:inside-the-earth:layers}{outer core}.

\textbf{12.} $5 \div 6.371 \approx 0.785$: the oceanic \omterm{def:g6:inside-the-earth:layers}{crust} must be
painted only \textbf{0.78~mm} thick on the 2-metre model, the thickness of
about eight sheets of paper ($0.78 \div 0.1 \approx 8$).
\end{solution}
